\section{The intrinsic nilpotent-depth filtration}
\label{sec:depth}

We isolate the functorial mechanism that unifies the two local models.
Let \(Z\) be any noetherian scheme and let \(\mathcal N_Z\) denote the
nilradical of \(\OO_Z\). Since \(Z\) is noetherian,
\(\mathcal N_Z\) is a coherent nilpotent ideal. For \(j\ge1\) define
\begin{equation}\label{eq:depth-filtration}
 \ZZ_j(Z)=
 V\!\left(\Ann_{\OO_Z}(\mathcal N_Z^{\,j})\right).
\end{equation}

\begin{lemma}[Functoriality]\label{lem:depth-functorial}
Every isomorphism \(f:Z\to Z'\) carries
\(\ZZ_j(Z)\) isomorphically to \(\ZZ_j(Z')\) for every \(j\).
The construction commutes with restriction to open subsets.
\end{lemma}
\begin{proof}
An isomorphism carries the nilradical to the nilradical, hence carries
its \(j\)-th power and its annihilator ideal to the corresponding
objects. Both formation of powers and annihilators commute with
localization.
\end{proof}

Apply this to \(Z=\widehat D_R\). On the intrinsic open obtained by
restricting to the smooth locus of \(Z_{\mathrm{red}}\), Proposition
\ref{prop:full-fitting} identifies \(Z\) with \(D_R\). Theorem
\ref{thm:normal-form} gives the local ring
\[
 A=\OO/(t(t,f))
\]
with nilradical generated by the class of \(t\). Therefore
\[
 \Ann_A((t))=(t,f)/(t(t,f)),
\]
so \eqref{eq:depth-filtration} recovers \(E_R\) scheme-theoretically.

\begin{proposition}[Depth one is the polarized ramification support]
\label{prop:depth-one}
On the smooth-reduction locus of \(\widehat D_R\),
\[
 \ZZ_1(\widehat D_R)=E_R,\qquad
 \widehat\NN_R|_{E_R}=\OO_{E_R}(-\widehat\Delta_R).
\]
Consequently the pair
\[
 \left(\ZZ_1(\widehat D_R),
 \omega_{\ZZ_1(\widehat D_R)}
 \otimes
 (\widehat\NN_R|_{\ZZ_1(\widehat D_R)})^{-9}\right)
\]
is intrinsic when \(e=4\), and its graded section ring is the fourth
Veronese of the polarized quartic section ring.
\end{proposition}
\begin{proof}
Only the last assertion is new notation. The annihilator calculation
above is Theorem~\ref{thm:normal-form}; the line-bundle identity and
section-ring reconstruction are Proposition~\ref{prop:canonical-line}
and Theorem~\ref{thm:main-reconstruction}.
\end{proof}

The point of \eqref{eq:depth-filtration} is that it makes no reference
to a projection chart. The closed schemes are automatically nested:
\[
 \widehat\NN_R^{\,j+1}\subset\widehat\NN_R^{\,j}
 \quad\Longrightarrow\quad
 \Ann(\widehat\NN_R^{\,j})
 \subset
 \Ann(\widehat\NN_R^{\,j+1}),
 \qquad
 \ZZ_{j+1}\subset\ZZ_j.
\]
The new issue is therefore not formal functoriality but the
scheme-theoretic computation of the layers as one crosses projection
rank. Theorem~\ref{thm:global-residual-colon} solves this globally:
on the whole Grassmannian the residual colon
\[
 \mathcal J_R^{\mathrm{res}}
 =
 (\mathcal I_{\widehat D_R}:\mathcal I_{\widehat\Delta_R})
\]
determines every annihilator and every associated-graded piece,
including all projection coranks and rank-drop loci. The remaining
task is geometric identification of those canonical colon schemes.

On the first higher-corank chart the next power is nonzero. The
calculation in \S\ref{sec:coranktwo} gives, on the simple-contact open,
\[
 \widehat\NN_R^3=0,\qquad
 \Ann(\widehat\NN_R^2)=\II_{\Sigma_R}/\II_{\widehat D_R}.
\]
Section~\ref{sec:ranktwo-boundary} removes the simple-root assumption.
Section~\ref{sec:stratified-nilpotent} computes the two layers
simultaneously across the corank-one/corank-two boundary. On the
mixed-regular neighbourhood,
\[
 \gr_{\widehat\NN_R}\OO_{\widehat D_R}
 =
 \OO_{\widehat\Delta_R}
 \oplus
 (\mathcal L\otimes\OO_{\mathcal C_R})
 \oplus
 (\mathcal L^{\otimes2}\otimes\OO_{\Sigma_R}),
\]
where \(\mathcal C_R\) is the unsplit contact cone and
\(\mathcal L\) is the Schubert conormal line. In particular
\(\ZZ_1\) specializes from the ramification support to its canonical
contact closure with the vertex thickening forced by the Fitting
algebra, while \(\ZZ_2\) is the rank-two vertex. On the first
tensor-rank boundary of the mixed map, \(\ZZ_2\) acquires a canonical
length-three nonreduced structure
(Theorem~\ref{thm:decomposable-kernel-wall}).

Thus
\[
 \ZZ_1\quad\leadsto\quad\ZZ_2
\]
is not merely a comparison of two calculations on disjoint strata:
it is the associated-graded specialization of one intrinsic
nilpotent algebra.

\begin{remark}[Relation with canonical filtrations on multiple structures]
\label{rem:multiple-filtrations}
Nilpotent and multiple structures have long been studied through
canonical filtrations and associated graded modules; ribbons are the
multiplicity-two prototype, and primitive multiple schemes carry a
canonical chain of lower multiplicities. See
\cite{BayerEisenbud,Manolache,Drezet}. The construction
\eqref{eq:depth-filtration} is deliberately more elementary: it is
available for the present Fitting schemes without assuming that they
are primitive, locally complete-intersection multiple structures, or
Cohen--Macaulay along every higher-corank stratum. The additional
content here is that consecutive annihilator layers and their
associated graded modules have a specific geometric meaning forced by
multiplication: the first carries the polarized ramification K3, the
second detects the rank-two Schubert support, and their specialization
through contact collisions is controlled by the unsplit contact
Cartier ideal rather than by a choice of primary roots.
\end{remark}
